\documentclass{handout}
\usepackage{handoutSetup}\usepackage{handoutShortcuts} 
\begin{document}
\handoutHeader

\begin{verbatimwrite}{\jobname.title}
The Prescott Real Business Cycle Model
\end{verbatimwrite}

\handoutNameMake

\begin{verbatimwrite}{./body.tex}

This handout presents the elements of the original
Real Business Cycle model of aggregate fluctuations, as laid out
by \cite{prescottTheoryAhead}, along with a few critiques articulated by \cite{summers:skeptical} and others.\opt{MarginNotes}{\marginpar{\tiny The good news is that you already understand the subtlest and hardest parts of the model - 
the intertemporal consumption optimization problem.
}}{}


Consider a representative household whose goal is to maximize
\begin{equation}
\Ex_{t} \left[\sum_{t=0}^{\infty} \Discount^{t}\uFunc(\cons_{t},\leisure_{t})      \right]
\end{equation}
where $\leisure_{t}$ is the fraction of time the representative agent spends at leisure (not working); the
alternative to leisure is the number of hours you work, which will be designated
$\labor_{t}$, and the time endowment is normalized to 1, so that
\begin{equation}\begin{gathered}\begin{aligned}
        \labor_{t}+\leisure_{t} & =  1.
\end{aligned}\end{gathered}\end{equation}

Assume that the structure of the utility function is
\begin{equation}\begin{gathered}\begin{aligned}
        \uFunc(\cons,\leisure) & =  \left(\frac{(\cons^{1-\zeta}\leisure^{\zeta})^{1-\CRRA}}{1-\CRRA}\right)
\\  \uFunc^{\cons}     & =  (\cons^{1-\zeta}\leisure^{\zeta})^{-\CRRA}\cons^{-\zeta}\leisure^{\zeta}(1-\zeta)
\\  \uFunc^{\leisure }  & =  (\cons^{1-\zeta}\leisure^{\zeta})^{-\CRRA}\cons^{1-\zeta}\leisure^{\zeta-1}\zeta.
\end{aligned}\end{gathered}\end{equation}

Think about the maximum amount of income that could be gained if the representative agent
worked every waking hour:\opt{MarginNotes}{\marginpar{\tiny Since the time
    endowment is one, resources are $\Wage_{t}*1$.}}{}
\begin{equation}\begin{gathered}\begin{aligned}
        \inc_{t} & =  \Wage_{t}.
\end{aligned}\end{gathered}\end{equation}

The representative agent can then think of deciding to `purchase' two things with this 
endowment of income: leisure $\leisure_{t}$ whose price is $\Wage_{t}$, or 
consumption, whose price is normalized to one.

Over the past century in the U.S., wages have risen very substantially, but hours 
worked have not declined much if at all. (\cite{rameyFrancisLeisure}).  What kind of utility 
function implies that the budget share of a good 
(leisure) remains constant even as the price of the good changes 
sharply?  A Cobb-Douglas utility function.  Hence the assumption that 
utility is obtained from a Cobb-Douglas aggregate of consumption and 
leisure is consistent with the lack of a strong trend in hours worked per 
worker.

Since workers are choosing how many hours to work
as well how much to consume, a first order condition
will characterize the optimal choice {\it between} consumption
and leisure {\it within} a period.  In particular, the price of leisure
is $\Wage_{t}$ and the price of consumption is 1, so the ratio of the
marginal utility of leisure to the marginal utility of consumption
should be 
\begin{equation}\begin{gathered}\begin{aligned}
        \left(\frac{\Wage_{t}}{1}\right) & =  \left(\frac{\uFunc^{\leisure }_{t}}{\uFunc^{\cons}_{t}}\right). \label{eq:focwucul}
\end{aligned}\end{gathered}\end{equation}

To see this, note that the consumer's goal is to
\begin{equation}
\max_{\{\cons_{t},\leisure_{t}\}} \uFunc(\cons_{t},\leisure_{t})     .
        \label{eq:intramax}
\end{equation}

Suppose the consumer has decided to spend a given amount $\chi_{t}$ in period $t$
on a combination of consumption and leisure,
\begin{equation}\begin{gathered}\begin{aligned}
        \cons_{t}+\Wage_{t}\leisure_{t} & =  \chi_{t}
\end{aligned}\end{gathered}\end{equation}

Then \eqref{eq:intramax} becomes
\begin{equation}
\max_{\{ \leisure_{t} \}} \uFunc(\chi_{t}-\Wage_{t}\leisure_{t},\leisure_{t})        
        \label{eq:intramax2}
\end{equation}
for which the FOC is
\begin{equation}\begin{gathered}\begin{aligned}
        - \uFunc^{\cons}_{t}\Wage_{t} + \uFunc_{t}^{\leisure } & =  0  
\\      \Wage_{t} & =  (\uFunc^{\leisure }_{t}/\uFunc_{t}^{\cons}).
\end{aligned}\end{gathered}\end{equation}

Returning to \eqref{eq:focwucul}
\begin{equation}\begin{gathered}\begin{aligned}
        \left(\frac{\Wage_{t}}{1}\right) & =  \left(\frac{\uFunc^{\leisure }_{t}}{\uFunc^{\cons}_{t}}\right)
\\ \Wage_{t} & =  \left(\frac{\cons_{t}^{1-\zeta}\leisure_{t}^{\zeta-1}}{\cons_{t}^{-\zeta}\leisure^{\zeta}_{t}}\right)\left(\frac{\zeta}{1-\zeta}\right)        
\\  & =  \left(\frac{\cons_{t}}{\leisure_{t}}\right)\left(\frac{\zeta}{1-\zeta}\right)   
\end{aligned}\end{gathered}\end{equation}
or
\begin{equation}\begin{gathered}\begin{aligned}
        \Wage_{t}\leisure_{t}/\cons_{t} & =  \left(\frac{\zeta}{1-\zeta}\right)      . \label{eq:wlc}
\end{aligned}\end{gathered}\end{equation}


Since we know that $\leisure $ has been roughly constant over long periods 
of time, this implies that as wages rise, consumption rises by roughly 
the same amount.

One of the original proimises of the DSGE literature was to calibrate
its business-cycle models based on either long-run facts (like the lack
of a trend in $\leisure$) or on micro data (like intertemporal elasticities
estimated using household data).  So how is $\zeta$ calibrated?

If wages are defined as per unit of labor, then if on average consumption
roughly equals labor income $\cons = (1-z)w$ we have
\begin{equation}\begin{gathered}\begin{aligned}
        \left(\frac{\wage  \leisure}{\wage (1-z)}\right) & =  \left(\frac{\zeta}{1-\zeta}\right)  \\
        \left(\frac{1-\leisure}{\leisure }\right) & =  \left(\frac{1-\zeta}{\zeta}\right)  
\\      1/\leisure - 1 & =  1/\zeta - 1
\\  \leisure & =  \zeta     .
\end{aligned}\end{gathered}\end{equation}

So $\zeta$ should be calibrated to be equal to the proportion of their 
available (i.e.\  non-sleep) time people spend not working.  A 40-hour 
work week (along with 8 hours of sleep a day) would yield $\zeta = 
2/3$.  Among other taste parameters, Prescott chooses log utility
($\lim_{\CRRA \rightarrow 1}\uFunc(\cons)$) and $\Discount = 0.96.$

The aggregate production function is assumed to be Cobb-Douglas,
\begin{equation}\begin{gathered}\begin{aligned}
        \Kap_{t+1} & =  \overbrace{\daleth}^{1-\depr} \Kap_{t}+\Inc_{t}-\Cons_{t}  \\
        \Inc_{t} & =  \PtyLev_{t}\Kap_{t}^{1-\labShare}\Labor_{t}^{\labShare}
\end{aligned}\end{gathered}\end{equation}
where $\labor_{t} = (1-\leisure_{t})\Hours_{t} = \labor_{t}\Hours_{t}$
where $\Hours_{t}=\sum_{i} \hours_{i}$ is the aggregate amount of
Hours available to members of the working population. Constant income
shares and perfect competition imply
\begin{equation}\begin{gathered}\begin{aligned}
        \FFunc_{\Labor}\Labor/\Inc & =  \labShare \PtyLev_{t}\Kap_{t}^{1-\labShare}\Labor^{\labShare-1}\Labor_{t}/\Inc_{t}  \\
         & =  \labShare
\end{aligned}\end{gathered}\end{equation}
so that labor's share of GDP is roughly constant.  Prescott sets 
labor's share to a constant 64 percent, and chooses a depreciation 
rate of $\depr = 0.10$.

The crucial assumption, however, is about the productivity process, 
since `technology shocks' are assumed to drive business cycles.

Prescott defines the `hat' operator $\hat{}$ as:
\begin{equation}\begin{gathered}\begin{aligned}
        \hat{X}_{t} & =  \left(\frac{X_{t}-X_{t-1}}{X_{t-1}}\right)
\end{aligned}\end{gathered}\end{equation}
which implies from the production function that
\begin{equation}\begin{gathered}\begin{aligned}
        \hat{\PtyLev}_{t} & \approx  \hat{Y}_{t}-\labShare \hat{\Labor}_{t} - (1-\labShare)\hat{\Kap}_{t}.
\end{aligned}\end{gathered}\end{equation}
(this is just the Solow residual).

Prescott `estimates' a productivity process that takes the form
\begin{equation}\begin{gathered}\begin{aligned}
        \hat{\PtyLev}_{t} & =  \ptyGro_{t} + \epsilon_{t}
\end{aligned}\end{gathered}\end{equation}
with a standard deviation of $\sigma_{\epsilon} = 0.76$ per quarter.  

Prescott makes sufficient assumptions (perfect competition, etc.)  so
that the social planner's problem is the same as the decentralized
solution.  With log utility, the social planner's problem is
\begin{equation}
        \max \Ex_{t}\left[\sum_{t=0}^{\infty}\Discount^{t} \left((1-\zeta) \log \cons_{t} + \zeta \log (1-\labor_{t})\right)\right] \label{eq:objlog}
\end{equation}
subject to
\begin{equation}\begin{gathered}\begin{aligned}
        \kap_{t+1} & =  \daleth \kap_{t}+\PtyLev_{t}k^{1-\labShare}_{t}\labor^{\labShare} - \cons_{t}  \\
        \hat{\PtyLev}_{t} & =  \ptyGro + \epsilon_{t}.
\end{aligned}\end{gathered}\end{equation}

Prescott argues that the way to judge the model is by whether it
produces plausible statistics for standard deviations of the 
key variables.  He produces a table that argues it does:

\begin{table}[ht]
\center
\begin{tabular}{lccc}
                 & $\sigma_{\leisure }$  & $\sigma_{y}$   $\sigma_{n}$ \\ \hline
   US Data       & 0.76          & 1.76           1.67   
\\ Model         & 0.76          & 1.48           0.76
\end{tabular}
\end{table}

Since the first column is calibrated, it isn't a test of the model.  
The second column comes out of the model, and isn't too bad a fit.  
However, the third column is a terrible fit.  What it says is that 
labor input is much more variable over the course of the business 
cycle than this model would suggest.  

What's going on?  To understand the answer, we need to understand why 
hours fluctuate in this model at all.  Recall that we deliberately constructed the 
model (by choosing a utility function that was Cobb-Douglas in consumption and 
leisure) in a way designed to {\it prohibit} any {\it long-run} response of hours 
worked to wages.  Since hours worked are being chosen freely on a
day-by-day basis by workers in this model, there must be some 
incentive that causes them to be willing to put up with {\it short-term}
variation in hours (over the business cycle).

The answer is that transitory productivity shocks provide an incentive 
to work harder some times than others.  In particular, if there is a 
temporary positive productivity shock you will be willing to work longer hours 
than usual, while if there is a negative productivity shock everybody 
wants to take a vacation.  \opt{MarginNotes}{\marginpar{\tiny Thus the 1930s in the US could be characterized as `The Great Vacation \ldots'.}}{}

To see this formally, consider again the first order conditions from 
the maximization problem.  We showed in \eqref{eq:wlc} that \opt{MarginNotes}{\marginpar{\tiny Intuition: u must be unchanged if $\leisure $ increases one and $\cons$ decreases by $\Wage_{t}$.}}{}
\begin{equation}\begin{gathered}\begin{aligned}
        \Wage_{t}\leisure_{t}/\cons_{t} & =  \left(\frac{\zeta}{1-\zeta}\right)  \\
         & =  \Wage_{t+1}\leisure_{t+1}/\cons_{t+1}.  \label{eq:wlcttp1}
\end{aligned}\end{gathered}\end{equation}

Now note that since \eqref{eq:objlog} is separable in consumption
and leisure the intertemporal FOC will imply that
\begin{equation}\begin{gathered}\begin{aligned}
        1/\cons_{t}   & =  \Rfree_{t+1}\Discount /\cons_{t+1}
\\  \cons_{t+1}   & =  \Rfree_{t+1}\Discount \cons_{t}. 
\end{aligned}\end{gathered}\end{equation}

Combining this with \eqref{eq:wlcttp1} gives
\begin{equation}\begin{gathered}\begin{aligned}
        \Wage_{t}\leisure_{t} & =  \Wage_{t+1}\leisure_{t+1}/\Rfree_{t+1}\Discount  \\
        \leisure_{t+1}/\leisure_{t} & =  \left(\frac{\Rfree_{t+1}\Discount}{\Wage_{t+1}/\Wage_{t}}\right)  \\
        \hat{\leisure }_{t+1} & \approx  -\hat{\wage }_{t+1}+(\rfree_{t+1}-\theta)
\end{aligned}\end{gathered}\end{equation}

What this equation tells us is that there are two ways to make $\leisure_{t}$
(and therefore $\labor_{t}$) fluctuate over the business cycle:
\begin{enumerate}
        \item  Make transitory movements in real wages induce a strong labor supply response
        \begin{itemize}
                \item  Problem:  A large number of microeconomic studies have estimated
the elasticity of labor supply with respect to wages to be much too small to explain
the observed fluctuations in hours over the business cycle.  Indeed, there is even controversy
about whether real wages rise or fall over the business cycle.  But even if we accept that 
wages fall during recessions, the evidence at the micro level does not seem to suggest that people are willing
to dramatically change their work hours in response to transitory fluctuations in wages
        \end{itemize}

        \item Make movements in interest rates induce a strong labor 
        supply response (the idea is that you will work hard during the 
        period of high interest rates in order to earn more cash 
        which can be invested to take advantage of the high interest rate - in thinking about this, 
        consider that a high value of $\Rfree_{t+1}$ will induce high leisure 
        growth between $t$ and $t+1$ by resulting in low leisure in period 
        $t$)
\begin{itemize}
        \item  Problem:  If this is the mechanism, there should at the same time be a strong consumption
        response:
        \begin{equation}\begin{gathered}\begin{aligned}
                \Wage_{t}\leisure_{t}/\cons_{t} & =  \Wage_{t+1}\leisure_{t+1}/\cons_{t+1}  \\
                \cons_{t+1}/\cons_{t} & =  \Wage_{t+1}\leisure_{t+1}/\Wage_{t}\leisure_{t}
        \end{aligned}\end{gathered}\end{equation}
     so if the labor supply response is not being driven by wage differences,
     there should be a one-for-one comovement of consumption with leisure - i.e.\
     recessions should be periods of high consumption and booms should be periods
     of low consumption!

\end{itemize}
\end{enumerate}

This latter is a quite general problem with the DSGE framework in which 
fluctuations in employment over the business cycle are driven by 
voluntary changes in hours worked.  

\renewcommand\refname{\noindent{\href{http://www.econ2.jhu.edu/people/ccarroll/public/lecturenotes/DSGEModels/RBC-Prescott/RBC-Prescott.bib}{\large References {\small (click to download \texttt{.bib} file)}}}}

\input handoutBibMake.tex

\end{verbatimwrite}
\input ./body.tex


\end{document}

